\section{Local graphs and fixed parameters}\label{sec:setup}

We first explain the class of graph integrals to be estimated.
Near each point of \(\Sigma\), choose tangent coordinates and a
normal in which the surface is a graph with positive definite
Hessian. At a boundary point, a smooth parametrization up to the
boundary extends to an open neighborhood of the parameter
half-disk; this is the usual smooth-boundary convention
\cite[Chapter~1, pp.~27--28]{Lee2013Smooth}. Projection onto the tangent plane has invertible
differential there. The inverse function theorem on the extension
therefore gives a graph on an open planar neighborhood, whose
restriction parametrizes the original surface piece.

After shrinking a chart, Lemma~\ref{lem:continuation} gives a fixed
invertible affine change of ambient frequency coordinates in which
the graph agrees locally with a phase \(\phi\in C^\infty(\R^2)\)
in the precise continuation class of Theorem~\ref{thm:packet}.
In particular, this phase satisfies
\begin{equation}\label{eq:ellipticity}
 cI\le D^2\phi\le CI,\qquad
 \sup_{\R^2}|D^j\phi|<\infty\quad(j\ge2).
\end{equation}
The elementary arguments below use only
Equation~\eqref{eq:ellipticity}; the packet estimate also uses the
continuation form verified in Lemma~\ref{lem:continuation}.

Fix a smooth nonnegative function \(\chi\) supported in a square
centered at zero, such that
\begin{equation}\label{eq:window-partition}
 \sum_{\nu\in a\Z^2}(\chi_a^\nu(\xi))^2=1,\qquad
 \chi_a^\nu(\xi)=\chi((\xi-\nu)/a),\qquad a>0.
\end{equation}
For example, start with a compactly supported bump positive on
\([-1,1]^2\), and divide it by the square root of the sum of the
squares of its integer translates. This denominator is smooth,
positive, and periodic. Let \(D_{\mathrm f}\ge1\) be a fixed side
length for a square containing the support of \(\chi\).

Write
\[
 E_\phi g(y)=\int_{\R^2}g(\xi)
       e\bigl(y'\cdot\xi+y_3\phi(\xi)\bigr)\dd\xi,
 \qquad y=(y',y_3)\in\R^2\times\R.
\]
A partition at scale one reduces each compactly supported graph
integral to finitely many functions of the form
\begin{equation}\label{eq:local-F}
 F(y)=E_\phi(h\chi_1^{\nu_*})(y).
\end{equation}
Here \(\nu_*\in\Z^2\) is fixed and \(h\in L^3(\R^2)\) is supported
in a fixed compact set \(Q_0\). To see the reduction, write graph
data \(g=\sum_{\nu_*}(g\chi_1^{\nu_*})\chi_1^{\nu_*}\), using
Equation~\eqref{eq:window-partition}; only finitely many terms meet
\(\supp g\). The extra copy of the window is included in \(h\).

All constants for a fixed setup may depend on \(\phi,Q_0,\chi\),
\(D_{\mathrm f}\), and \(\nu_*\). These objects, and all label ranges
used later, are chosen before the large scale or the small loss
exponents vary. In particular,
\begin{equation}\label{eq:velocity-monotonicity}
 (\nabla\phi(\xi)-\nabla\phi(\eta))\cdot(\xi-\eta)
 \ge c|\xi-\eta|^2
\end{equation}
on \(\R^2\). Thus the velocity map \(U_\xi=-\nabla\phi(\xi)\) is
bi-Lipschitz on each fixed bounded region and has bounded size there.

The final graph estimate controls Equation~\eqref{eq:local-F} by
\(\|h\|_3\). We record the precise surface Jacobians in
Section~\ref{sec:surfaces}.
